% GENERATED FILE. Edit canonical lessons, then run npm run generate:books.
% canonical-source-hash: fnv1a64:def6b07ed31c5362
% canonical-lessons: PA-C115-rajai, PA-C115-cullha, PA-C115-bhada, PA-C115-kaici, PA-C115-sui

\chapter{Quilt, Stove, Vessel, Scissors, Needle}
\label{ch:pa-quilt-stove-vessel-scissors-needle}
\hlchaptermodality{115}

\begin{chapteropening}
\textbf{By the end of this chapter:} \emph{I can say a quilt, a stove, a vessel, scissors and a needle.}

Complete the last lesson of chapter 115: I can say a quilt, a stove, a vessel, scissors and a needle.
\end{chapteropening}

\section[\texorpdfstring{\pa{ਰਜਾਈ} (rajāī)}{ਰਜਾਈ (rajāī)}]{\pa{ਰਜਾਈ} (rajāī) — a quilt}

\label{lesson:PA-C115-rajai}



\begin{quote}
Before the new one: say the Punjabi for a bedsheet, then the Punjabi for a blanket.
\end{quote}

\subsection*{You'll want to know: \pa{ਰਜਾਈ}}
\textbf{\pa{ਰਜਾਈ}} — \emph{rajāī} — \textquotedblleft{}a quilt\textquotedblright{}.

\begin{grammarlens}[title={What you've built}]
One more word for this chapter.
\end{grammarlens}

\subsection*{Your turn}
\begin{itemize}
\raggedright
  \item \emph{Say it:} \emph{rajāī}
  \item \emph{Say it:} \emph{rajāī}, once more
  \item \emph{Say it:} say \emph{rajāī}
  \item \emph{Recall:} say the Punjabi for a bedsheet, then the Punjabi for a blanket, then say \emph{rajāī} again
\end{itemize}

\begin{tcolorbox}[breakable,colback=teal!4,colframe=teal!35!black,title={Before you move on}]
What does \pa{ਰਜਾਈ} mean? (\textquotedblleft{}A quilt\textquotedblright{}.) Say it once more.
\end{tcolorbox}
\section[\texorpdfstring{\pa{ਚੁੱਲ੍ਹਾ} (cullhā)}{ਚੁੱਲ੍ਹਾ (cullhā)}]{\pa{ਚੁੱਲ੍ਹਾ} (cullhā) — a stove}

\label{lesson:PA-C115-cullha}



\begin{quote}
Before the new one: say the Punjabi for a blanket, then the Punjabi for a quilt.
\end{quote}

\subsection*{You'll want to know: \pa{ਚੁੱਲ੍ਹਾ}}
\textbf{\pa{ਚੁੱਲ੍ਹਾ}} — \emph{cullhā} — \textquotedblleft{}a stove\textquotedblright{}.

\begin{grammarlens}[title={What you've built}]
One more word for this chapter.
\end{grammarlens}

\subsection*{Your turn}
\begin{itemize}
\raggedright
  \item \emph{Say it:} \emph{cullhā}
  \item \emph{Say it:} \emph{cullhā}, once more
  \item \emph{Say it:} say \emph{cullhā}
  \item \emph{Recall:} say the Punjabi for a blanket, then the Punjabi for a quilt, then say \emph{cullhā} again
\end{itemize}

\begin{tcolorbox}[breakable,colback=teal!4,colframe=teal!35!black,title={Before you move on}]
What does \pa{ਚੁੱਲ੍ਹਾ} mean? (\textquotedblleft{}A stove\textquotedblright{}.) Say it once more.
\end{tcolorbox}
\section[\texorpdfstring{\pa{ਭਾਂਡਾ} (bhā̃ḍā)}{ਭਾਂਡਾ (bhā̃ḍā)}]{\pa{ਭਾਂਡਾ} (bhā̃ḍā) — a vessel, a utensil}

\label{lesson:PA-C115-bhada}



\begin{quote}
Before the new one: say the Punjabi for a quilt, then the Punjabi for a stove.
\end{quote}

\subsection*{You'll want to know: \pa{ਭਾਂਡਾ}}
\textbf{\pa{ਭਾਂਡਾ}} — \emph{bhā̃ḍā} — \textquotedblleft{}a vessel, a utensil\textquotedblright{}.

\begin{grammarlens}[title={What you've built}]
One more word for this chapter.
\end{grammarlens}

\subsection*{Your turn}
\begin{itemize}
\raggedright
  \item \emph{Say it:} \emph{bhā̃ḍā}
  \item \emph{Say it:} \emph{bhā̃ḍā}, once more
  \item \emph{Say it:} say \emph{bhā̃ḍā}
  \item \emph{Recall:} say the Punjabi for a quilt, then the Punjabi for a stove, then say \emph{bhā̃ḍā} again
\end{itemize}

\begin{tcolorbox}[breakable,colback=teal!4,colframe=teal!35!black,title={Before you move on}]
What does \pa{ਭਾਂਡਾ} mean? (\textquotedblleft{}A vessel, a utensil\textquotedblright{}.) Say it once more.
\end{tcolorbox}
\section[\texorpdfstring{\pa{ਕੈਂਚੀ} (kaĩcī)}{ਕੈਂਚੀ (kaĩcī)}]{\pa{ਕੈਂਚੀ} (kaĩcī) — scissors}

\label{lesson:PA-C115-kaici}



\begin{quote}
Before the new one: say the Punjabi for a stove, then the Punjabi for a vessel.
\end{quote}

\subsection*{You'll want to know: \pa{ਕੈਂਚੀ}}
\textbf{\pa{ਕੈਂਚੀ}} — \emph{kaĩcī} — \textquotedblleft{}scissors\textquotedblright{}.

\begin{grammarlens}[title={What you've built}]
One more word for this chapter.
\end{grammarlens}

\subsection*{Your turn}
\begin{itemize}
\raggedright
  \item \emph{Say it:} \emph{kaĩcī}
  \item \emph{Say it:} \emph{kaĩcī}, once more
  \item \emph{Say it:} say \emph{kaĩcī}
  \item \emph{Recall:} say the Punjabi for a stove, then the Punjabi for a vessel, then say \emph{kaĩcī} again
\end{itemize}

\begin{tcolorbox}[breakable,colback=teal!4,colframe=teal!35!black,title={Before you move on}]
What does \pa{ਕੈਂਚੀ} mean? (\textquotedblleft{}Scissors\textquotedblright{}.) Say it once more.
\end{tcolorbox}
\section[\texorpdfstring{\pa{ਸੂਈ} (sūī)}{ਸੂਈ (sūī)}]{\pa{ਸੂਈ} (sūī) — a needle}

\label{lesson:PA-C115-sui}



\begin{quote}
Before the new one: say the Punjabi for a vessel, then the Punjabi for scissors.
\end{quote}

\subsection*{You'll want to know: \pa{ਸੂਈ}}
\textbf{\pa{ਸੂਈ}} — \emph{sūī} — \textquotedblleft{}a needle\textquotedblright{}.

\begin{grammarlens}[title={What you've built}]
One more word for this chapter.
\end{grammarlens}

\subsection*{Your turn}
\begin{itemize}
\raggedright
  \item \emph{Say it:} \emph{sūī}
  \item \emph{Say it:} \emph{sūī}, once more
  \item \emph{Say it:} say \emph{sūī}
  \item \emph{Recall:} say the Punjabi for a vessel, then the Punjabi for scissors, then say \emph{sūī} again
\end{itemize}

\begin{tcolorbox}[breakable,colback=teal!4,colframe=teal!35!black,title={Before you move on}]
What does \pa{ਸੂਈ} mean? (\textquotedblleft{}A needle\textquotedblright{}.) Say it once more.
\end{tcolorbox}
